% -------------------------------------------------------------------------
% WBS ID: RES-ML-ALG
% Deliverable: Learning Algorithm and Architecture Assessment
% Status: In progress
% Evidence status: Algorithm assessment deferred by current hydrodynamic scope
% Review status: Not reviewed
% Last updated: 2026-07-19
% Dependencies: RES-ML-DAT, RES-ML-REP, RES-ML-TRN
% Downstream interfaces: RES-ML-SUR, RES-ML-SPEC
% -------------------------------------------------------------------------

\section{RES-ML-ALG --- Learning Algorithm and Architecture Assessment}
\label{sec:res-ml-alg}

\subsection{Purpose and Scope}
\label{sec:res-ml-alg-purpose}

This Deliverable records future algorithm candidates only. The active
project does not assess or select an ML algorithm; it first produces
validated hydrodynamic outputs and structured provenance.

\subsection{Research Questions}
\label{sec:res-ml-alg-research-questions}

% TODO: State the questions that must be answered before the Deliverable
% can be accepted.

\subsection{Terminology and Definitions}
\label{sec:res-ml-alg-definitions}

% TODO: Define the technical terminology, symbols, quantities and
% classifications used in this Deliverable.

\subsection{Literature Evidence}
\label{sec:res-ml-alg-literature-evidence}

% TODO: Record evidence from peer-reviewed papers, authoritative datasets,
% standards, technical reports and primary documentation.
%
% Do not create a detached literature-summary list. Explain how each source
% supports, contradicts or limits a project decision.

The Intelligent Systems lecture material may be used as terminology and
foundation context for linear regression, polynomial regression, logistic
regression, Naive Bayes, model selection, kernels, support-vector methods,
clustering, dimensionality reduction, neural networks, linear algebra, calculus,
probability and statistics. It shall not be used as the sole evidence for
research claims that require peer-reviewed or authoritative primary sources.

\subsection{Governing Relations and Quantitative Definitions}
\label{sec:res-ml-alg-governing-relations}

% TODO: Record relevant equations, dimensionless groups, empirical models,
% extraction rules or quantitative definitions.
%
% Use align environments for displayed equations.
% Every displayed equation must have a unique \label{eq:...}.
% Do not place punctuation after displayed equations.

\subsection{Machine-Learning Role}
\label{sec:res-ml-alg-machine-learning-role}

% TODO: Complete this RES-ML Domain-specific subsection.

\subsection{Non-ML Baseline}
\label{sec:res-ml-alg-non-ml-baseline}

% TODO: Complete this RES-ML Domain-specific subsection.

\subsection{Input and Output Representation}
\label{sec:res-ml-alg-input-output-representation}

% TODO: Complete this RES-ML Domain-specific subsection.

\subsection{Dataset Requirements}
\label{sec:res-ml-alg-dataset-requirements}

% TODO: Complete this RES-ML Domain-specific subsection.

\subsection{Model or Architecture Candidates}
\label{sec:res-ml-alg-model-architecture-candidates}

If ML is reactivated later, the assessment hierarchy may be:
\begin{enumerate}
  \item linear regression baseline;
  \item regularised polynomial regression;
  \item logistic regression for classification;
  \item linear and nonlinear support-vector methods;
  \item kernel-based regression;
  \item tree ensembles or boosting;
  \item Gaussian-process methods where computationally feasible;
  \item feed-forward neural networks only when dataset volume supports them;
  \item PCA or another controlled dimensionality-reduction method for
    correlated outputs or histories;
  \item clustering for exploratory analysis only.
\end{enumerate}

\subsection{Training and Validation Method}
\label{sec:res-ml-alg-training-validation-method}

% TODO: Complete this RES-ML Domain-specific subsection.

\subsection{Evaluation Metrics}
\label{sec:res-ml-alg-evaluation-metrics}

% TODO: Complete this RES-ML Domain-specific subsection.

\subsection{Generalisation and Uncertainty}
\label{sec:res-ml-alg-generalisation-uncertainty}

% TODO: Complete this RES-ML Domain-specific subsection.

\subsection{Simulation and Optimisation Integration}
\label{sec:res-ml-alg-simulation-optimisation-integration}

% TODO: Complete this RES-ML Domain-specific subsection.

\subsection{Candidate Approaches}
\label{sec:res-ml-alg-candidate-approaches}

% TODO: Define the credible candidate approaches considered.

\subsection{Critical Comparison}
\label{sec:res-ml-alg-critical-comparison}

Model selection shall consider dataset size, dimensionality, nonlinearity,
uncertainty, interpretability, extrapolation, failure-boundary behaviour and
implementation cost. Popularity or implementation convenience alone is
insufficient evidence for selection.

\subsection{Assumptions and Applicability}
\label{sec:res-ml-alg-assumptions}

% TODO: State assumptions and the conditions under which they are valid.

\subsection{Limitations and Failure Conditions}
\label{sec:res-ml-alg-limitations}

% TODO: State where the evidence, model or selected approach ceases to be
% reliable or applicable.

\subsection{Project-Specific Conclusions}
\label{sec:res-ml-alg-conclusions}

The initial algorithm strategy is baseline-first and uncertainty-aware. A final
algorithm remains unresolved until data volume, feature dimensionality and
validation behaviour are known.

\subsection{Selected Approach and Rationale}
\label{sec:res-ml-alg-selected-approach}

% TODO: Record the selected approach only after sufficient evidence exists.
% State the alternatives rejected and the reasons for rejection.

\subsection{Unresolved Decisions}
\label{sec:res-ml-alg-unresolved-decisions}

% TODO: Record unresolved questions, required evidence and decision owners.

\subsection{Requirements and Handoffs}
\label{sec:res-ml-alg-requirements}

% TODO: State requirements passed to other Research Domains, Software,
% verification activities or later execution work.

\subsection{Deliverable Completion Record}
\label{sec:res-ml-alg-completion-record}

% TODO: Record whether the Deliverable is:
%
% - Not started
% - In progress
% - Draft complete
% - Reviewed
% - Accepted
%
% Acceptance requires:
%
% - the research questions to be answered;
% - claims to be supported by appropriate evidence;
% - assumptions and limitations to be explicit;
% - the project-specific conclusion to be stated;
% - downstream requirements to be traceable;
% - unresolved decisions to be closed or formally deferred.
